\documentclass[12pt,a4paper]{article}

\usepackage[margin=1in]{geometry}
\usepackage{graphicx}
\usepackage{float}
\usepackage{longtable}
\usepackage{array}
\usepackage{booktabs}
\usepackage{hyperref}
\usepackage{xcolor}
\usepackage{listings}
\usepackage{titlesec}

\titleformat{\section}{\large\bfseries}{\thesection}{1em}{}
\titleformat{\subsection}{\normalsize\bfseries}{\thesubsection}{1em}{}

\lstset{
    basicstyle=\ttfamily\small,
    breaklines=true,
    frame=single,
    columns=fullflexible
}

\title{Mini Project: Image Enhancement for Poor Quality Images}
\author{Student Name \\ ID}
\date{\today}

\begin{document}

\maketitle
\tableofcontents
\newpage

\section{Introduction}
This project applies different image enhancement techniques to selected poor-quality images. 
Each image is processed using methods suitable for its degradation type, such as blur, noise, low contrast, or poor visual clarity.

\section{Objectives}
The objectives of this mini project are:
\begin{itemize}
    \item Extract important components from selected images
    \item Enhance blurry images
    \item Remove noise from noisy images
    \item Improve the visual quality of challenging images
\end{itemize}

\section{Tools and Libraries}
The project was implemented in Python using:
\begin{itemize}
    \item OpenCV
    \item NumPy
    \item Matplotlib
    \item scikit-image
    \item SciPy
\end{itemize}

\section{Methodology}
For each image:
\begin{enumerate}
    \item The original image was inspected visually
    \item More than one suitable enhancement method was applied
    \item Intermediate outputs were saved
    \item The best final output was selected
    \item Results were discussed in terms of pros, cons, limitations, and possible future work
\end{enumerate}

\section{Task 1: Component Extraction}

\subsection{7circles}
\textbf{Techniques used:}
\begin{itemize}
    \item Grayscale conversion
    \item Gaussian blurring
    \item Otsu thresholding
    \item Morphological opening and closing
    \item Circle/component extraction
\end{itemize}

\textbf{Why these techniques were selected:} \\
These techniques are suitable for isolating circular shapes from the background while reducing noise and improving contour quality.

\textbf{Code:} \\
See \texttt{src/componentExtraction.py}

\textbf{Discussion:} \\
Pros: Good isolation of circular components. \\
Cons: May miss faint circles if contrast is weak. \\
Future work: Hough Circle Transform tuning or adaptive preprocessing.

\subsection{COVID-19 Chart}
\textbf{Techniques used:}
\begin{itemize}
    \item Grayscale conversion
    \item CLAHE
    \item Adaptive thresholding
    \item Morphological operations
    \item Contour extraction
\end{itemize}

\textbf{Why these techniques were selected:} \\
Charts contain thin lines, labels, and structured elements, so local contrast enhancement and threshold refinement help separate components.

\textbf{Code:} \\
See \texttt{src/componentExtraction.py}

\textbf{Discussion:} \\
Pros: Improves component visibility and boundary extraction. \\
Cons: Some fine text may merge with other regions. \\
Future work: Color-based segmentation in HSV if chart colors are distinct.

\section{Task 2: Blur Enhancement}

\subsection{building}
\textbf{Techniques used:}
\begin{itemize}
    \item Unsharp masking
    \item Laplacian sharpening
    \item Frequency-domain Wiener-like sharpening
\end{itemize}

\textbf{Why these techniques were selected:} \\
The building image contains structural edges and lines, making sharpening methods effective for visual restoration.

\textbf{Discussion:} \\
Pros: Better edge visibility and stronger structures. \\
Cons: Aggressive sharpening may create halos. \\
Future work: Blind deconvolution with estimated blur kernels.

\subsection{dog}
\textbf{Techniques used:}
\begin{itemize}
    \item Mild unsharp masking
    \item Bilateral filtering followed by sharpening
    \item Laplacian sharpening
\end{itemize}

\textbf{Why these techniques were selected:} \\
The dog image likely contains soft textures, so enhancement must preserve natural appearance without over-amplifying noise.

\textbf{Discussion:} \\
Pros: Better fur and edge clarity. \\
Cons: Strong sharpening can make the image unnatural. \\
Future work: More advanced deblurring methods with noise-aware restoration.

\section{Task 3: Noise Removal}

\subsection{text}
\textbf{Techniques used:}
\begin{itemize}
    \item Median filtering
    \item Adaptive thresholding
    \item Morphological cleanup
\end{itemize}

\textbf{Why these techniques were selected:} \\
Text images require strong noise suppression while preserving letter strokes and readability.

\textbf{Discussion:} \\
Pros: Good text/background separation. \\
Cons: Tiny characters may become thinner or broken. \\
Future work: Stroke-aware morphology or OCR-oriented enhancement.

\subsection{rocket}
\textbf{Techniques used:}
\begin{itemize}
    \item Median filtering
    \item Bilateral filtering
    \item Non-local means denoising
\end{itemize}

\textbf{Why these techniques were selected:} \\
These methods reduce noise while preserving rocket boundaries and major structures.

\textbf{Discussion:} \\
Pros: Better denoising with preserved edges. \\
Cons: Over-smoothing may remove fine details. \\
Future work: Estimate noise type more precisely before choosing the filter.

\subsection{windChart}
\textbf{Techniques used:}
\begin{itemize}
    \item Median filtering
    \item Bilateral filtering
    \item CLAHE
    \item Otsu thresholding
\end{itemize}

\textbf{Why these techniques were selected:} \\
The image contains chart data and labels, so denoising must preserve lines and readability.

\textbf{Discussion:} \\
Pros: Better readability and cleaner chart structure. \\
Cons: Thresholding may remove some weak lines. \\
Future work: Line-specific enhancement or adaptive local denoising.

\section{Task 4: Visual Enhancement}

\subsection{newsPaper}
\textbf{Techniques used:}
\begin{itemize}
    \item Grayscale conversion
    \item CLAHE
    \item Adaptive thresholding
    \item Morphological cleanup
\end{itemize}

\textbf{Why these techniques were selected:} \\
Newspaper images often suffer from low contrast and non-uniform illumination, so local contrast enhancement is effective.

\textbf{Discussion:} \\
Pros: Improved text readability and page structure. \\
Cons: Noise may be amplified in very weak areas. \\
Future work: Denoising before contrast enhancement.

\subsection{namePlate}
\textbf{Techniques used:}
\begin{itemize}
    \item Gamma correction
    \item CLAHE
    \item Sharpening
    \item Adaptive thresholding
\end{itemize}

\textbf{Why these techniques were selected:} \\
The image may suffer from poor lighting and weak character visibility, making gamma and local contrast enhancement appropriate.

\textbf{Discussion:} \\
Pros: Better plate detail visibility. \\
Cons: Thresholding may not always help if reflections are strong. \\
Future work: Reflection suppression and local illumination correction.

\section{Conclusion}
Different images require different enhancement strategies. 
A single method is not suitable for all degradations, so the project applied task-specific techniques and compared intermediate outputs to select the best final result.

\section{Future Work}
Possible future improvements include:
\begin{itemize}
    \item Blind deconvolution for blur restoration
    \item Better noise estimation for adaptive denoising
    \item Color-space based chart extraction
    \item OCR-guided enhancement for text-heavy images
\end{itemize}

\end{document}